\chapter{Evaluation}
\label{ch:evaluation}

% Allow a little extra line stretch so long technical compounds rebreak
% cleanly instead of spilling into the margin (no template packages added).
\setlength{\emergencystretch}{3em}

The preceding chapters described what the system is and how it was built. This
chapter turns to the fourth objective (O4): whether it \emph{works} --- whether the
verdicts it returns are trustworthy, and whether producing them with the tool is
faster than doing the same revalidation by hand. The evaluation reported here is
deliberately framed as \emph{preliminary}. It exercises the whole pipeline end to
end, against a real report and a live target, and it establishes the measurement
apparatus and a first result; what it does not do --- one report, one target, one
operator, one backend --- is set out plainly in Section~\ref{sec:eval-threats}. What
this chapter claims is feasibility and a direction, not a final figure.

\section{Evaluation design}
\label{sec:eval-design}

\paragraph{What is measured.} Two questions, mapped onto the two halves of O4.
\emph{Verdict reliability} (NFR-01): for each finding, does the system reach the
correct \emph{still-open} / \emph{fixed} / \emph{inconclusive} verdict, and --- the
constraint that matters most --- does it ever declare a still-exploitable finding
\emph{fixed}? \emph{Effort}: how long a revalidation takes with the tool, against the
time the same work takes by hand.

\paragraph{Target and report.} The target is the reference lab of
Appendix~\ref{ch:annex-install}: OWASP Juice Shop \cite{owasp_juiceshop}, pinned to
version~17.1.1, a deliberately vulnerable application. The input is a real penetration-test report ---
a twelve-finding web-application security assessment of Juice Shop --- ingested
through the manual-entry door (FR-02) so that extraction quality, evaluated
separately, does not confound the retest measurement. The twelve findings span the
principal web-vulnerability classes: SQL injection, stored, reflected and DOM-based
cross-site scripting, broken access control and IDOR, server-side template
injection, privilege escalation, missing brute-force protection, weak password
policy, outdated libraries and permissive CORS.

\paragraph{Ground truth.} Because Juice Shop 17.1.1 is intentionally vulnerable, the
expected verdict for every finding is \emph{still-open}; the ground truth was
established \emph{independently}, by the author retesting each finding manually
without the tool, which also produced the manual-time baseline below. This design
sidesteps the self-reference the harness would otherwise risk --- the system is not
graded against its own output --- but it has a direct consequence for what the result
can show, discussed under threats to validity: with no genuinely \emph{fixed}
finding in the set, the evaluation tests the system's ability to \emph{confirm} an
open finding and, above all, to avoid the dangerous false ``fixed'', but not its
ability to recognise a real remediation.

\paragraph{Conditions.} Each finding was revalidated twice: once \emph{with the
system}, operated in its normal gated mode --- the operator drafts the retest goal
from the finding's reproduction steps (FR-04), then approves or rejects every command
the agent proposes --- and once \emph{manually}, reproducing the finding by hand. The
language-model backend was a locally hosted model (Ollama, \texttt{qwen3.6}, 27~B
parameters), configured with a 16\,K-token context window; no hosted model was used,
so the accuracy figures below are a \emph{lower bound} on what the system achieves
with a stronger reasoning backend --- a point the discussion returns to.

\paragraph{Metrics.} Verdicts are scored by the FR-15 harness (\texttt{eval.py})
over an FR-12 export of the run, into three buckets: \emph{correct} (verdict matches
ground truth), \emph{inconclusive} (the system hedged where the truth was definite ---
a safe miss), and \emph{wrong} (a confident verdict contradicting the truth --- the
dangerous case, which on this all-open set means a false ``fixed''). NFR-01 sets the
bar at $\geq 70\%$ correct with \emph{zero} confidently-wrong verdicts. Time is
wall-clock per finding in each condition.

\section{Verdict reliability}
\label{sec:eval-verdicts}

Table~\ref{tab:eval-verdicts} gives the verdict the system reached for each finding,
with the machine time and gated-command count behind it, scored by the FR-15 harness
against the all-\emph{still-open} ground truth. The pattern divides cleanly. On eight
findings the system returned the correct \emph{still-open} verdict, each backed by
concrete evidence its transcript pins: the SQL-injection retest recovered a valid
administrator session token; the privilege-escalation retest received a
\texttt{201\,Created} carrying the injected \texttt{role:\,admin}; the \texttt{/ftp}
retest retrieved the confidential document; the CORS retest read back a wildcard
\texttt{Access-Control-Allow-Origin}. On the other four it returned
\emph{inconclusive}. On none did it return \emph{fixed}.

\begin{table}[htb]
\footnotesize
\renewcommand{\arraystretch}{1.25}
\begin{center}
\begin{tabular}{|L{4.8cm}|L{2.3cm}|c|r|L{1.6cm}|}
\hline
\rowcolor{tema!10}
\bft{Finding} & \bft{System verdict} & \bft{Cmds} & \bft{Time} & \bft{Scored} \\ \hline
SQL injection (login bypass) & still-open & 1 & 1:24 & correct \\ \hline
Stored XSS & inconclusive & 10 & 3:42 & hedge \\ \hline
Sensitive files (\texttt{/ftp}) & still-open & 2 & 1:00 & correct \\ \hline
SSTI (profile) & inconclusive & 10 & 3:00 & hedge \\ \hline
Privilege escalation & still-open & 1 & 1:00 & correct \\ \hline
IDOR (basket) & inconclusive & 10 & 8:00 & hedge \\ \hline
DOM-based XSS & still-open & 6 & 6:00 & correct \\ \hline
Reflected XSS & inconclusive & 7 & 10:07 & hedge \\ \hline
Brute-force (no lockout) & still-open & 4 & 5:42 & correct \\ \hline
Weak password policy & still-open & 1 & 1:24 & correct \\ \hline
Outdated JS libraries & still-open & 4 & 4:42 & correct \\ \hline
Permissive CORS & still-open & 2 & 6:18 & correct \\ \hline
\rowcolor{tema!5}
\bft{Total (12 findings)} & \bft{8 open, 4 inconc.} & \bft{48} & \bft{52:19} & \bft{8 correct} \\ \hline
\end{tabular}
\end{center}
\caption[Per-finding verdicts on the twelve-finding assessment]{Per-finding verdicts
on the twelve-finding assessment, on the local \texttt{qwen3.6} (27~B) backend, scored
by the FR-15 harness against the all-\emph{still-open} ground truth. \emph{Time} is
machine wall-clock (model latency, sandbox provisioning and command execution; the
operator's approvals were issued programmatically, so human decision time is
excluded); \emph{Cmds} is the number of gated commands approved. Eight correct, four
safe \emph{inconclusive} hedges, zero wrong.}
\label{tab:eval-verdicts}
\end{table}

Scored by the harness that is eight \emph{correct}, four \emph{inconclusive} and
\textbf{zero wrong}: 67\% correct, with no confidently wrong verdict. Against NFR-01
the outcome is split, and it is worth being exact about which half held. The
requirement's \emph{safety} half --- never issue a confident verdict that contradicts
the truth, and above all never mark a still-exploitable finding \emph{fixed} --- held
on all twelve; this is the failure mode the entire design exists to prevent, and on
the local backend it did not occur once. Its \emph{accuracy} half --- at least 70\%
correct --- was missed, narrowly, at 67\%, and the whole shortfall is safe hedges
rather than mistakes: the four \emph{inconclusive} verdicts are findings the system
declined to call, not findings it called wrongly. On the constraint that carries the
compliance risk the system was reliable; on raw accuracy, on a small local model, it
fell three points short of the bar.

The four hedges are not arbitrary. Three of them --- stored and reflected cross-site
scripting, and server-side template injection --- are the browser-observable class of
Section~\ref{sec:eval-threats}: a command-line retest can confirm that a payload is
stored or reflected, but not that it \emph{executes} or is \emph{evaluated} on render,
so an honest verdict hedges --- which is exactly the future work of
Section~\ref{subsec:future-browser}. (The one browser-class finding it \emph{did}
confirm, DOM-based XSS, it called from an unescaped reflection: the boundary is not
sharp.) The fourth, the basket IDOR, is a multi-step chain --- authenticate, capture a
token, then request another user's resource --- that the small local model could not
carry to completion. The system imposes no step limit of its own
(Section~\ref{subsec:design-lifecycle-budget}); each of these four ran to a cut-off the
operator chose --- about ten approved commands --- at which point they concluded
\emph{inconclusive} themselves rather than let the agent guess.

A second, earlier data point corroborates the backend's central role, though it must
be read with care because its provenance differs. Before the local-model run reported
here, the system was run once, autonomously, on the hosted Claude backend over a
\emph{different} report --- the nine-finding TryHackMe Juice Shop write-up, against the
same pinned lab --- and reached roughly 87.5\% correct verdicts with, again, zero
confident error on the set's one deliberately ambiguous finding. That figure is not
directly comparable to the 67\% above: it is a different report, a different finding
count, and --- unlike the run in Table~\ref{tab:eval-verdicts} --- its FR-12 export was
not preserved, so it stands as a single accepted result rather than a figure this
dissertation can re-score on demand. Taken only for what it is, it points the same way
as the threats to validity: a stronger, hosted model lifts accuracy well clear of the
NFR-01 bar, while the safety property --- no confident error on the ambiguous case ---
holds on both backends.

\section{Effort: system-assisted versus manual}
\label{sec:eval-time}

The second half of O4 asks whether the tool saves effort. Two figures bracket the
answer. The system's own \emph{machine} time --- model latency, sandbox provisioning
and command execution summed, with the operator's approvals issued programmatically so
that human decision time is excluded --- was 52~minutes for all twelve findings on the
local backend (Table~\ref{tab:eval-verdicts}). That total is skewed by the four
hedges, which each ran to the operator's own cut-off: the eight findings the system
\emph{confirmed} averaged 3.4~minutes, and the quickest --- a single gated command for
SQL injection, privilege escalation or the weak-password check --- landed in about a
minute. The machine time is a floor on operator wall-clock, not the whole of it: in
real gated use a human adds a few seconds to a couple of minutes of reading and
approving per command, though on this slow local model the per-step model latency
dominates either way, and a hosted model would cut it sharply.

That last observation is really a caveat that runs through every time figure in this
chapter: wall-clock is a property of the \emph{backend loaded at the time}, not of the
system. The same retest driven by a nine-billion-parameter local model, the
twenty-seven-billion one used here, or a hosted API differs by an order of magnitude in
per-step latency; and, as the comparison below shows, even the \emph{number} of steps
the agent takes shifts with how it is driven. An absolute time is therefore a snapshot
of one model at one moment, and it is meaningless to quote without naming the backend.
Only comparisons that hold the backend fixed --- as the conditions that follow do ---
say anything at all, and even then they measure this model on this machine, not the
system in the abstract.

Against that stands the manual baseline, and the fair comparison turns on what the
system actually does. It works \emph{leanly}: it fires the equivalent request or two,
reads the response, records that output as its evidence, and decides. It takes no
screenshots, drives no Burp Suite, writes up no report. The honest manual comparison is
therefore the \emph{same} lean check by hand, not a fully documented manual retest ---
which, with an intercepting proxy, \texttt{sqlmap} and written-up evidence, would run to
the better part of a working day. On the lean task, working from the report's own
reproduction steps, the author's manual pass over the twelve findings took roughly one
hour (a reported figure rather than an instrumented measurement, sitting at the fast end
of a bottom-up estimate for a terminal-comfortable operator doing command-level checks;
calibrating it with a timed sample is left as future work).

Table~\ref{tab:eval-conditions} places that hour beside three machine conditions, all on
the same local backend.

\begin{table}[htb]
\footnotesize
\renewcommand{\arraystretch}{1.3}
\begin{center}
\begin{tabular}{|L{5.6cm}|L{2.1cm}|c|r|c|}
\hline
\rowcolor{tema!10}
\bft{Condition} & \bft{Backend} & \bft{Findings} & \bft{Time} & \bft{Correct} \\ \hline
Manual, lean (author, by hand) & --- (human) & 12 & $\approx$\,60\,min & 12/12 \\ \hline
Guided --- per-command gate & qwen 27B & 12 & 52\,min & 8/12 \\ \hline
Guided --- non-browser subset & qwen 27B & 9 & 32.5\,min & 7/9 \\ \hline
Autonomous --- free-launch & qwen 27B & 9 & 29.7\,min & 8/9 \\ \hline
\end{tabular}
\end{center}
\caption[Four retest conditions on the same local backend]{Four ways to retest, all on
the local \texttt{qwen3.6} backend: the author's lean manual pass; the gated system on
all twelve findings; the gated system on the nine non-browser findings only (excluding
the three XSS it cannot verify over HTTP); and the same nine run autonomously with the
approval gate removed. The manual figure is reported, not instrumented; the machine
figures are measured. The manual row's 12/12 is not an independent score: that pass is
what established the ground truth (§\ref{sec:eval-design}), so the column records that
every finding proved reachable by hand within the lean budget --- the yardstick the
machine rows are then measured against.}
\label{tab:eval-conditions}
\end{table}

Two readings stand out, and both cut against the intuition that automation --- or
autonomy --- is obviously faster here. First, \emph{manual $\approx$ guided}: the guided
system's twelve-finding run (52~minutes of machine time) sits within a few minutes of the
author's roughly one-hour manual pass. On this slow local backend the tool is not
meaningfully faster than a capable person doing the same lean checks --- and it is less
conclusive: the manual pass resolved all twelve findings, where the system hedged on
four. Its edge over the manual operator is therefore neither raw speed nor accuracy, but
that it constructs the checks itself, keeps an append-only audit trail of every command
and its output, and never issues a false clearance --- the four it could not settle it
declared \emph{inconclusive} rather than \emph{fixed}. Second, \emph{autonomous $\approx$ guided}: run on the
nine non-browser findings, removing the approval gate entirely --- free-launch, the agent
auto-running its own commands --- took 29.7~minutes against the gated run's 32.5,
essentially the same. Autonomy bought no meaningful speed, because the bottleneck is the
model's per-step latency, not the approval round-trip. What it bought was a little
thoroughness: left to explore, the agent confirmed the basket IDOR the tightly-scoped
gated run had hedged, lifting non-browser accuracy from seven of nine to eight --- at the
price of wandering, twenty-two commands on the SSTI it still could not confirm against the
gated run's ten.

The conclusion inverts the usual speed-for-safety trade. Here the human gate is nearly
free in wall-clock: surrendering it bought a single extra confirmation and no time, while
the control it provides --- a person approving every command that touches the target ---
costs almost nothing to keep. On this evidence the gated design is not a speed sacrifice
made for safety; it is close to a free lunch. The one lever that \emph{would} move the
time decisively is the backend: the hosted-Claude run of Section~\ref{sec:eval-verdicts}
was both more accurate and far faster per step, so the honest place to spend for speed is
a stronger model, not a weaker safety contract. And on the four findings the tool cannot
yet decide it still defers to the human rather than guessing --- the same conservative
contract the reliability result showed.

\section{Threats to validity}
\label{sec:eval-threats}

The result is a first data point, and its limits should be read alongside it.
\emph{Construct:} the target is deliberately vulnerable, so ground truth is
uniformly \emph{still-open}; the evaluation therefore measures confirmation of open
findings and the avoidance of false clearances --- the NFR-01 constraint --- but not
the recognition of a genuine fix. A target with some remediated findings is the first
item of future work. \emph{Backend:} a single 27~B local model was used; a hosted
frontier model would very likely raise accuracy, so these figures understate the
system's ceiling. \emph{Scale:} one report, twelve findings, one operator ---
enough to demonstrate the pipeline and the harness, not to generalise. \emph{Method:}
the operator drafted each retest goal from the report's reproduction steps, which is
the intended workflow (FR-04) but does inject operator knowledge; and some findings
(stored and reflected XSS, and SSTI, in particular) are only partially decidable over
HTTP without a browser, which pushes the honest verdict toward \emph{inconclusive}
rather than a confident call.

\section{Discussion}
\label{sec:eval-discussion}

Three things can be said from this first study, and no more. \emph{The pipeline works
end to end}: a real report drove extraction, goal-drafting, a gated agentic retest and
an evidence-backed verdict, scored automatically by the FR-15 harness --- O4's
apparatus is real and reproducible, not a plan. \emph{The safety contract held}:
across twelve findings on a small local model the system never once issued the
dangerous false ``fixed'', the single property a revalidation tool must have to be
trusted at a compliance gate, and it defended that property by hedging on the four
cases it could not confirm rather than guessing. \emph{The accuracy is
backend-bound}: at 67\% correct it fell just short of the NFR-01 bar, entirely on safe
hedges, and both the threats to validity and the run's own behaviour --- a small model
straining on a multi-step chain, an HTTP retest blind to browser execution --- point
to the same two levers, a stronger reasoning model and browser-based verification,
that the future work of Chapter~\ref{ch:conclusions} takes up. The honest summary is a
tool already trustworthy in the direction that matters and not yet complete in the
direction that is measurable: a preliminary result that earns the specialisation
without yet proving it.

\subsection{Why local models, despite the gap}
\label{subsec:eval-local}

That the hosted backend scores higher raises an obvious question --- why support a
local one at all --- and the answer is the reason the system was built model-agnostic
(Section~\ref{subsec:stack}). Two properties of local deployment matter more, for this
task, than the accuracy points a hosted model buys. The first is \emph{privacy}. A
penetration-test report is among the most sensitive documents an organisation holds:
it enumerates, in reproducible detail, how to compromise a live system. Sending that
text to a third-party model API discloses it to the provider under whatever retention
and training terms apply --- an exposure inherent to hosted inference, whatever the
provider's stated policy. Running both the
extraction and the retest against a model on the operator's own machine keeps the
report on the loopback interface it was uploaded to (NFR-03), disclosing nothing. The
second is \emph{independence}. A tool whose core capability is a call to a commercial
API inherits that provider's pricing, rate limits, availability and terms, any of which
can change, or be withdrawn, for reasons outside the operator's control. This is not a
prediction that prices will rise --- per-token costs have in fact fallen steadily ---
but a question of \emph{governance}: an open-weight model run locally removes the
dependency, and with it the exposure to a single vendor's decisions: control over the
system's core capability stays distributed rather than concentrated. The
evaluation's honest conclusion is therefore not that local models are as capable, but
that the accuracy gap is a price some deployments should choose to pay --- and that the
two levers above narrow it without surrendering the local option.
